% !TEX root = ../nc-quadric-classification.tex
% 2026-11-21 / v1.35
\section{相対 Picard 作用による cubic Serre 平方根の構成}
\label{sec:relative-picard-roots}

\paragraph{主張の範囲と依存.}
本稿は標数零における\emph{非線型} geometric cubic 五つ組の普遍族上に、
étale 局所に実際の二周期を構成する。
その平方は四次 tensor の巡回から定まる canonical 四周期に
頂点スカラーを除いて一致する。
全ての恒等式を被約な普遍基底上で証明してから引き戻すので、
非被約な Artin 基底も結論に含む。
原典の既知入力は、Artin--Tate--Van den Bergh の
\emph{相対} Picard 作用と、Van den Bergh の geometric 五つ組の構成である。
相対 tails と有限 foundation の比較、および Serre の双加群比較には、
定理~\ref{thm:relative-cubic-tails-tilting} と定理~\ref{thm:bimodule-line-comparison} を用いる。
本稿は任意の二周期が Serre 平方根であるとは主張しない。
相対化の新規性は未確認である。

\subsection{一次資料の正確な内容}

Van den Bergh \cite[Proposition 5.6.1]{VdB2011} は、
代数閉体 $k$、$\operatorname{char}k\ne2$ 上の全ての quadric
$\mathbb Z$ 代数 $D$ に対して $D\simeq D(2)$ を示す。
全 $\mathbb Z$ 代数の二周期性が主張されている。
その後の Corollary 5.6.2 が偶数部分の通常の次数付き環による表示を与える。

非線型の場合、楕円 helix $(C,(\mathcal L_i)_i)$ に対し
\[
 N=\mathcal L_2\mathcal L_0^{-1}\in\Pic^0(C),\qquad
 \mathcal L_{i+2}\simeq\mathcal L_iN
\]
である。$a\in\Pic^0(C)$、$a^{-2}=N$ を選び、
\cite[Theorem 4.2.3]{VdB2011} の作用
\[
 \eta:\Pic^0(C)\longrightarrow\Aut(C),\qquad
 \eta(a)^*L=L\otimes a^{-\deg L}
 \tag{PR-1}\label{eq:field-eta}
\]
を用いる。$\sigma=\eta(a)$ は各次数二の束 $\mathcal L_i$ を $\mathcal L_{i+2}$ へ移す。
原典の Proposition 6.6 は二周期を標準 hull へ持ち上げるが、
任意の二周期 $\beta$ に対する $\beta^2=\theta$ を主張してはいない。

\begin{proposition}[ATV の相対 Picard 作用]\label{prop:ATV-relative-action}
$C\to S$ を $\PP^2_S$ の三次因子、または
$\PP^1_S\times_S\PP^1_S$ の bidegree $(2,2)$ の因子の平坦族とする。
$P=\Pic^0_{C/S}$ を、各幾何ファイバーの各被約既約成分で次数零の
直線束の類を表す相対 Picard 部分群とする。
$P$ の $C$ への作用があり、その構成は任意の底変更および族の同型に自然である。
作用の符号を \eqref{eq:field-eta} と合わせて $\eta$ と書く。
\end{proposition}

\begin{proof}
完全な証明は相対モジュライ論文 \cite[\texttt{relative_picard_roots}]{QuadricModuli2026} に置く。
\end{proof}

\subsection{被約な普遍基底で証明してから底変更する}

以下、$k$ は代数閉体、$\operatorname{char}k=0$ とする。
四つの二次元空間 $V_0,\ldots,V_3$ を固定し、
\[
 \mathscr G\subset\PP(V_0\otimes V_1\otimes V_2\otimes V_3)
\]
を geometric 五つ組の開部分とする。
$02\mid13$ flattening の階数一の閉部分を除き、
\[
 S=\mathscr G\setminus\mathscr G_{\mathrm{lin}}
\]
と置く。$S$ は滑らかで被約、既約な有限型 $k$ スキームである。
以下の構成は $S$ の affine open 上で行って貼り合わせてもよい。

普遍直線 $W\subset(V_0\otimes\cdots\otimes V_3)\otimes\OO_S$ を巡回し、
$V_{i+4}=V_i$ として $W_i$ と三次関係 $R_i$ を定める。
これにより普遍 cubic $\mathbb Z$ 代数 $D$ を得る。
$W_i$ は全て同じ基底線 $W$ の巡回像であり、この共通の同定の下で
canonical 四周期
\[
 \theta:D\longrightarrow D(4)
 \tag{PR-2}\label{eq:theta}
\]
は各 $V_i=V_{i+4}$ 上の恒等である。
この記述に $W$ の大域的な生成元は必要ない。
局所生成元 $w_i$ を使えば
\[
 w_i=\sum_a v_a\otimes r_a,
 \qquad w_{i+1}=\sum_a r_a\otimes\theta_i(v_a)
\]
という \cite[\S5.2]{VdB2011} の定義そのものである。

\begin{lemma}[普遍族の次数別平坦性]\label{lem:universal-flatness}
各 $D_{ij}$ は有限局所自由で、標準 cubic 解消は各内部次数で分裂完全である。
その双対複体も AS 双対の所定の次数以外で完全で、残る cohomology は可逆層である。
従ってこれらの提示・解消・双対・関係の交わりは任意の底変更と両立する。
\end{lemma}

\begin{proof}
完全な証明は相対モジュライ論文 \cite[\texttt{relative_picard_roots}]{QuadricModuli2026} に置く。
\end{proof}

普遍 point curve を $\pi:C\to S$ とする。
隣接二因子の評価による $2\times2$ 行列の determinant が
$\PP(V_0^\vee)\times\PP(V_1^\vee)$ 内に $C$ を定める。
geometric 条件は、determinant が零である点でもこの行列の階数が一であることを保証する。
三点 correspondence の隣接二点への射影は、単元小行列式による掃き出しから
スキームとして同型になる。
従って全 $\mathcal L_i$ を同じ $C$ 上の直線束として定義できる。
これは第~\ref{sec:relative-cubic-noetherian}節の相対 point-curve 構成である。

$S$ は線型 locus を除いているので、determinant は各幾何ファイバーで非零である。
したがって $C/S$ は平坦な relative Cartier divisor である。
その標準 exact sequence
\[
 0\longrightarrow\OO_{\PP^1_S\times\PP^1_S}(-2,-2)\otimes\pi^*\ell
 \longrightarrow\OO_{\PP^1_S\times\PP^1_S}
 \longrightarrow\OO_C\longrightarrow0
 \tag{PR-3}\label{eq:curve-sequence}
\]
（$\ell$ は基底の可逆層）から、$\pi_*\OO_C=\OO_S$ が普遍的に成り立ち、
$R^1\pi_*\OO_C$ は可逆層、$R^{>1}\pi_*\OO_C=0$ となる。

\begin{lemma}[被約基底での相対 Picard の一致判定]\label{lem:reduced-equalizer}
$X$ を被約スキーム、$Y\to X$ を相対 diagonal が immersion である代数空間とする。
二つの section $f,g:X\to Y$ が全幾何点上で一致すれば、$f=g$ である。
この判定は $Y=\Pic_{C/X}$ に適用できる。
\end{lemma}

\begin{proof}
完全な証明は相対モジュライ論文 \cite[\texttt{relative_picard_roots}]{QuadricModuli2026} に置く。
\end{proof}

\begin{lemma}[普遍 helix の相対的な差]\label{lem:relative-helix-difference}
$N=[\mathcal L_2\mathcal L_0^{-1}]$ は $P=\Pic^0_{C/S}$ の section であり、相対 Picard 群で
\[
 [\mathcal L_{i+2}\mathcal L_i^{-1}]=N
 \tag{PR-4}\label{eq:helix-difference}
\]
が全ての $i$ について成り立つ。
\end{lemma}

\begin{proof}
完全な証明は相対モジュライ論文 \cite[\texttt{relative_picard_roots}]{QuadricModuli2026} に置く。
\end{proof}

\subsection{平方根の étale パラメータ空間}

\begin{lemma}[相対 Picard の二等分]\label{lem:picard-halves}
$P=\Pic^0_{C/S}$ は滑らかな相対次元一の可換群代数空間であり、
\[
 [-2]:P\longrightarrow P
\]
は étale かつ全射である。
従って
\[
 T=S\mathbin{\times}_{N,P,[-2]}P
 \tag{PR-5}\label{eq:T}
\]
は $S$ 上 étale 全射であり、普遍 section $a\in P(T)$ は
$a^{-2}=N$、すなわち $a^2=N^{-1}$ を満たす。
\end{lemma}

\begin{proof}
完全な証明は相対モジュライ論文 \cite[\texttt{relative_picard_roots}]{QuadricModuli2026} に置く。
\end{proof}

$T$ は $k$ 上滑らかで被約である。
必要なら $T$ の étale scheme chart を取り、以下では $T$ 自体をスキームとして扱う。
命題~\ref{prop:ATV-relative-action} により、
普遍 Picard class $a$ から実際の曲線の自己同型
\[
 \sigma=\eta(a):C_T\longrightarrow C_T
 \tag{PR-6}\label{eq:sigma}
\]
が得られる。
ここでは $a$ を表す大域的な直線束を選ぶ必要はない。
群作用そのものが相対 Picard 空間上に定義されている。

\begin{lemma}[線束の比較と二周期]\label{lem:relative-beta}
$T$ の Zariski 開被覆上で同型
\[
 \varepsilon_i:\sigma^*\mathcal L_i\xrightarrow{\sim}\mathcal L_{i+2}
 \qquad(i=0,1,2,3)
 \tag{PR-7}\label{eq:epsilons}
\]
を選べ、それらが実際の代数同型
\[
 \beta:D_T\xrightarrow{\sim}D_T(2)
 \tag{PR-8}\label{eq:beta}
\]
を与える。
\end{lemma}

\begin{proof}
完全な証明は相対モジュライ論文 \cite[\texttt{relative_picard_roots}]{QuadricModuli2026} に置く。
\end{proof}

\subsection{平方と canonical 四周期の比較}

\begin{theorem}[相対的な適合 Serre 根]\label{thm:relative-compatible-root}
上の $T$ をさらに Zariski 局所化した基底上で、
\eqref{eq:beta} は頂点対角型の自己同型 $\delta$ に対し
\[
 \beta^2=\delta\theta
 \tag{PR-9}\label{eq:beta-square}
\]
を満たす。
従って tails 上では $(2)_\beta^2\simeq(4)_\theta$ である。
この同型は基底上に構成されており、任意の底変更後にも成り立つ。
特に非被約基底へ引き戻しても、相対 Serre 関手を持つ有限 foundation の比較の下で
\[
 s=(2)_\beta,\qquad s^2\simeq S^{-1}[2]
 \tag{PR-10}\label{eq:Serre-square}
\]
を得る。
\end{theorem}

\begin{proof}
完全な証明は相対モジュライ論文 \cite[\texttt{relative_picard_roots}]{QuadricModuli2026} に置く。
\end{proof}



\begin{remark}[比較線束の局所自明化]
相対 Picard の等式からは、まず二つの直線束の差が基底から来ることが分かる。
\eqref{eq:epsilons} の前に四つの比較直線束を自明化したことが必要である。
その選択を変えると $\beta_i$ は基底単元倍されるが、
projective な平方比較は変わらず、(d) の頂点スカラーがその差を吸収する。
さらに (e) では Serre 双加群の比較線束 $\ell$ を自明化している。
従って幾何ファイバー上の平方根の存在だけから、
基底線束を省略した大域的な平方同型を主張してはいない。
必要なのは有限個の基底直線束の局所自明化である。
\end{remark}

\begin{remark}[四項 foundation と右加群の向き]
Abelian helix から $D_{ij}=\Hom(E_i,E_j)$ を回復する場合は、
右 tails の構成に添字を反転した反対代数
\[
 C_{ij}=D_{-j,-i},\qquad \beta_C=\beta_D^{-1}
\]
を用いる。
すると $E_i$ は $\pi e_{-i}C$ に対応し、四項窓の自己準同型代数は
元の foundation の積と一致する。
$\beta_D^2=\delta\theta_D$ を反転すれば、$C$ 上でも
$\beta_C^2=\theta_C$ が頂点スカラーを除いて成り立つ。
ここで $\theta_C$ は $\theta_D^{-1}$ に対応する canonical 四周期である。
(e) の双加群 $P,Q$ とその比較 locus は、この $C$ の右 tails 上で直接取れる。
従って Serre の比較と右軌道積の規約は同じ向きで成立する。
\end{remark}

\subsection{Artin 持上げと atlas 構成への入力}

\begin{corollary}[有限型の étale 根 chart]\label{cor:finite-type-root-chart}
$S$ を覆う有限個の affine scheme $T_\alpha$ と étale 射 $T_\alpha\to S$ を選び、
その各々で上の実際の二周期および Serre 平方同型を構成できる。
従って $\coprod_\alpha T_\alpha$ は有限型の滑らかな $k$ スキームであり、
全ての非線型 cubic geometric 五つ組の点を覆う。
\end{corollary}

\begin{proof}
完全な証明は相対モジュライ論文 \cite[\texttt{relative_picard_roots}]{QuadricModuli2026} に置く。
\end{proof}

\begin{corollary}[指定した Picard 根の全 Artin 持上げ]\label{cor:all-artin-root}
$x\in S(k)$ と $a_x\in\Pic^0(C_x)$、$a_x^{-2}=N_x$ を固定する。
\eqref{eq:T} の中の $a_x$ に対応する点を含む affine étale scheme chart を選び、
その点を $t$ とする。必要なら前の有限被覆にこの chart を一枚追加する。
任意の局所 Artin $k$-代数 $A$ と、その閉点が $x$ となる framed geometric
五つ組 $\operatorname{Spec} A\to S$ に対し、$t$ を指定した $\operatorname{Spec} A\to T$ の持上げは一意に存在する。
対応する二周期は Serre 適合であり、任意の Artin 商・小拡大に整合する。
\end{corollary}

\begin{proof}
完全な証明は相対モジュライ論文 \cite[\texttt{relative_picard_roots}]{QuadricModuli2026} に置く。
\end{proof}

\paragraph{Abelian 変形と atlas への適用.}
相対根の構成自体は、非線型 locus 全体で特異・可約・非被約な曲線ファイバーを含む。
平坦 Abelian Artin 変形からの適用では、
補題~\ref{lem:cubic-artin-overlap-recovery} と命題~\ref{prop:framed-cubic-artin-lift} が
DDL による helix の持上げから、可逆部分束 $W_0$ とその overlap・縮約を含む
指定基底付き五つ組を回復する。
補題~\ref{lem:artin-two-heart-comparison} は、同じ四項 foundation を持つ二つの圏の
標準 dg Morita 同値が、指定された exact な還元を保って heart-exact となることを示す。
従って系~\ref{cor:all-artin-root} の根を元の Abelian 変形へ移せる。
この指定同型の持上げと有限型効果性を合わせた smooth atlas は
定理~\ref{thm:quintuple-chart-artin-lifting} と定理~\ref{thm:cubic-smooth-open} で構成する。
線型 cubic 点は全て pointed $(\operatorname{coh}(\PP^1\times\PP^1),\OO)$ に対応し、
既存の可換二次曲面の étale Morita-root chart を用いればよい。

